\documentclass[11pt]{article}
\usepackage{amsmath,amssymb,amsthm}
\usepackage[margin=1.2in]{geometry}
\usepackage{microtype}
\usepackage{parskip}
\usepackage{xcolor}
\usepackage{booktabs}

\setlength{\parindent}{0pt}
\setlength{\parskip}{6pt}

\begin{document}
\pagestyle{empty}

\begin{center}
  {\large\bfseries RetroDiff: Training Objective Derivation}
\end{center}

\medskip

\noindent
Let $x$ denote a clean data sample and $x_t$ its noised version at diffusion step $t \in \{1,\ldots,T\}$.
The forward process $q(x_{1:T}|x)$ gradually corrupts $x$; the model $p_\theta$ learns to reverse it.
Applying Jensen's inequality to $\log p_\theta(x)$ yields the variational lower bound (ELBO):

\begin{align}
  \log p_\theta(x)
    &\geq \mathbb{E}_{q(x_{1:T}|x)}\!\left[\log\frac{p_\theta(x_{0:T})}{q(x_{1:T}|x)}\right]
     = -\mathcal{L}_{\mathrm{ELBO}}
  \label{eq:elbo}
\end{align}

\noindent
The bound in \eqref{eq:elbo} is tight when the variational posterior matches the true reverse process.
Expanding the joint $p_\theta(x_{0:T})$ via the Markov factorisation and grouping terms by diffusion step
decomposes $\mathcal{L}_{\mathrm{ELBO}}$ into a reconstruction term and a sum of per-step KL divergences:

\begin{align}
  \mathcal{L}_{\mathrm{ELBO}}
    &= \underbrace{\mathbb{E}_q[-\log p_\theta(x_0|x_1)]}_{\text{reconstruction}}
     + \sum_{t=2}^{T}
       \underbrace{\mathrm{KL}\!\left(q(x_{t-1}|x_t,x) \,\|\, p_\theta(x_{t-1}|x_t)\right)}_{\text{step-}t\text{ denoising}}
  \label{eq:decomp}
\end{align}

\noindent
Each KL term in \eqref{eq:decomp} measures how well the learned reverse step $p_\theta(x_{t-1}|x_t)$
approximates the tractable posterior $q(x_{t-1}|x_t,x)$.
\textbf{RetroDiff}'s key contribution is to condition every denoising step on a set
$R(x) = \{r_1,\ldots,r_k\}$ of $k$ exemplars retrieved from the training corpus based on $x$.
The per-sample ELBO is evaluated under this conditioning, and the final training objective
averages over the data distribution:

\begin{align}
  \mathcal{L}_{\mathrm{RetroDiff}}(\theta)
    &= \mathbb{E}_{x \sim p_{\mathrm{data}}}\!\left[
         \mathcal{L}_{\mathrm{ELBO}}\!\left(x \,\big|\, R(x)\right)
       \right]
  \label{eq:retro}
\end{align}

\noindent
Minimising \eqref{eq:retro} encourages the model to leverage retrieved exemplars as soft structural
priors at every denoising step, without modifying the diffusion schedule or the ELBO derivation itself.

\end{document}
